\documentclass[10pt,aspectratio=169]{beamer}
\usepackage[T2A]{fontenc}
\usepackage[utf8]{inputenc}
\usepackage[russian]{babel}
\usepackage{amsmath,amssymb}
\usepackage{graphicx}
\usepackage{hyperref}
\usepackage{booktabs}
\usepackage{listings}

\usetheme{Madrid}
\usecolortheme{default}

\title{Последовательные методы в Simulation‑Based Inference}
\subtitle{Применимость Sequential методов к FMPE и дрейф эмбеддинга}
\date{\today}

\begin{document}

\frame{\titlepage}

\begin{frame}{Содержание}
\tableofcontents
\end{frame}

\section{Обзор последовательных SNPE методов}

\begin{frame}{Семейство Sequential Neural Posterior Estimation}
Классические методы коррекции смещения, возникающего при итеративном сужении proposal:
\begin{itemize}
    \item \textbf{SNPE‑A} (Papamakarios \& Murray, 2016) – точная коррекция для MoG.
    \item \textbf{SNPE‑B} (Lueckmann et al., 2017) – взвешивание по важности.
    \item \textbf{SNPE‑C / APT} (Greenberg et al., 2019) – контрастивная потеря (InfoNCE).
    \item \textbf{Truncated SNPE} (Deistler et al., 2022) – усечение proposal по априору.
\end{itemize}
\pause
Общая идея: на раунде $t$ сэмплируем $\theta \sim \tilde{p}_t(\theta)$, генерируем $x$ и обучаем $q_\phi(\theta|x)$, избегая смещения.
\end{frame}

\begin{frame}{Что требуется от модели $q_\phi(\theta|x)$}
Во всех SNPE нужно уметь:
\begin{itemize}
    \item Сэмплировать $\theta \sim q_\phi(\theta|x_o)$ (proposal).
    \item \textbf{Но также считать} $\log q_\phi(\theta|x)$:
    \begin{itemize}
        \item В SNPE‑A/B – для коррекции смещения/весов.
        \item В SNPE‑C (APT) – для вычисления контрастивной потери InfoNCE.
    \end{itemize}
\end{itemize}
\pause
Таким образом, ключевое требование к параметрической модели – быстрый доступ к логарифму плотности.
\end{frame}

\section{Особенности FMPE}

\begin{frame}{Flow Matching Posterior Estimation}
FMPE (Dax et al., 2023) – генеративная модель, основанная на Flow Matching.
\begin{itemize}
    \item Позволяет сэмплировать из $q_\phi(\theta|x)$ очень быстро (прямое интегрирование ODE).
    \item Плотность $\log q_\phi(\theta|x)$ \textbf{существует}, но её вычисление требует обратного интегрирования ODE с расчётом следа якобиана – \textbf{дорого}.
\end{itemize}
\pause
Обучение FMPE идёт через Conditional Flow Matching (CFM) – нужны только сэмплы и векторное поле, \textbf{логарифм плотности не используется}.
\end{frame}

\section{Применимость SNPE к FMPE}

\begin{frame}{Какие методы не подходят}
\begin{itemize}
    \item \textbf{SNPE‑A} и \textbf{SNPE‑B}: требуют $\log \tilde{p}(\theta)$ (плотности proposal) или $\log q_\phi(\theta|x)$ – для FMPE это вычислительный bottleneck.
    \item \textbf{SNPE‑C (APT)}: вопреки интуиции, \textbf{тоже требует} $\log q_\phi(\theta|x)$ в потере InfoNCE:
    \[
    \mathcal{L}_{\text{InfoNCE}} = -\frac{1}{K}\sum_{k=1}^K \log \frac{\exp(\log q_\phi(\theta_k|x_k))}{\sum_j \exp(\log q_\phi(\theta_j|x_k))}
    \]
    Без быстрого логарифма обучение становится неэффективным.
\end{itemize}
\pause
Вывод: прямое применение классических SNPE к FMPE невозможно без потери ключевого преимущества.
\end{frame}

\begin{frame}{Truncated Sequential NPE – единственный простой вариант}
Метод Truncated (усечённый) не использует никаких явных плотностей:
\begin{itemize}
    \item Proposal $\tilde{p}_t(\theta)$ обрезается до границ априора.
    \item Смещение отсутствует по построению, достаточно стандартной CFM‑потери.
\end{itemize}
\pause
\textbf{Truncated FMPE}:
\begin{enumerate}
    \item Обучить FMPE на данных из априора (раунд 0).
    \item Построить proposal = усечённый FMPE(·|$x_o$).
    \item Сэмплировать из этого proposal, симулировать $x$, добавить в обучающую выборку.
    \item Обучить FMPE заново или дообучить (с CFM‑лоссом) на всех данных.
    \item Повторить с шага 2 с новым proposal.
\end{enumerate}
Никаких $\log q_\phi$ не требуется.
\end{frame}

\begin{frame}{Перспективы других sequential подходов}
Авторы FMPE упоминают, что sequential NPE к FMPE «применим, но оставлен на будущее». Вероятные пути:
\begin{itemize}
    \item Использовать отдельный лёгкий суррогат для оценки $\log q_\phi$ (но вырастет сложность).
    \item Перейти к sequential Neural Ratio Estimation (не требует плотности модели).
    \item Разработать варианты APT с аппроксимацией $\log q_\phi$ без полного интегрирования ODE.
\end{itemize}
Пока Truncated FMPE – самый практичный выбор.
\end{frame}

\section{Дрейф эмбеддинга в sequential NPE}

\begin{frame}{Проблема изменяющегося представления $x_o$}
Если совместно с $q_\phi$ обучается summary‑сеть $h_\psi$:
\begin{itemize}
    \item На раунде $t$ используется представление $h_{\psi_t}(x_o)$.
    \item При переходе к раунду $t+1$ веса $\psi$ обновляются, и $h_{\psi_{t+1}}(x_o) \neq h_{\psi_t}(x_o)$.
    \item Proposal, основанный на старом представлении, становится несогласованным с новым.
\end{itemize}
\pause
Это нарушает предположения sequential методов и может привести к смещению или расходимости.
\end{frame}

\begin{frame}{Как решается в библиотеке \texttt{sbi}}
Стратегия по умолчанию:
\begin{enumerate}
    \item На первом раунде summary‑сеть обучается вместе с плотностью.
    \item Начиная со второго раунда, \textbf{веса summary‑сети замораживаются}.
    \item Пользователь может явно передать embedding\_net в density\_estimator, и \texttt{sbi} автоматически исключит её из оптимизатора после первого раунда.
\end{enumerate}
\pause
Другие подходы:
\begin{itemize}
    \item Предобучение эмбеддинга и полная заморозка с самого начала.
    \item Использование отдельного внешнего цикла для дообучения эмбеддинга (редко).
\end{itemize}
\end{frame}


\end{document}